% GENERATED FILE. Edit canonical lessons, then run npm run generate:books.
% canonical-source-hash: fnv1a64:847adfcfeec2df25
% canonical-lessons: HI-C430-divala, HI-C430-nidhi, HI-C430-puji, HI-C430-cuknaintr, HI-C430-ugahna

\chapter{Capital and Bankruptcy}
\label{ch:hi-capital-and-bankruptcy}
\hlchaptermodality{430}

\begin{chapteropening}
\textbf{By the end of this chapter:} \emph{I can talk about capital and funds, collecting money, a debt paid off and a business gone bankrupt.}

Complete the last lesson of chapter 430: I can talk about capital and funds, collecting money, a debt paid off and a business gone bankrupt.
\end{chapteropening}

\section[\texorpdfstring{\hi{दिवाला} (divālā)}{दिवाला (divālā)}]{\hi{दिवाला} (divālā) — bankruptcy}

\label{lesson:HI-C430-divala}



\begin{quote}
Before the new one: say the Hindi for a currency, then the Hindi for to cash.
\end{quote}

\subsection*{You'll want to know: \hi{दिवाला}}
\textbf{\hi{दिवाला}} — \emph{divālā} — \textquotedblleft{}bankruptcy\textquotedblright{}.

\begin{quote}
\hi{उस} \hi{दुकानदार} \hi{का} \hi{दिवाला} \hi{निकल} \hi{गया।}
\end{quote}

\begin{quote}
\emph{us dukāndār kā divālā nikal gayā.}
\end{quote}

\begin{quote}
That shopkeeper went bankrupt.
\end{quote}

Masculine.

\hi{दिवाला} \hi{निकलना}, to go bankrupt.

\begin{cousinweb}
\textbf{\hi{दिवाला}} (\emph{divālā}) is filed by CDIAL under Sanskrit \emph{dīpāvali}, \textquotedblleft{}a row of lamps\textquotedblright{} (CDIAL 6358), beside \hi{दिवाली}; how a row of lamps came to mean bankruptcy is not explained there.
\end{cousinweb}

\begin{grammarlens}[title={What you've built}]
One more word for getting around.
\end{grammarlens}

\subsection*{Your turn}
\begin{itemize}
\raggedright
  \item \emph{Say it:} \emph{divālā}
  \item \emph{Say it:} \emph{divālā}, once more
  \item \emph{Say it:} \emph{us dukāndār kā divālā nikal gayā.} — That shopkeeper went bankrupt.
  \item \emph{Recall:} say the Hindi for a currency, then the Hindi for to cash, then say \emph{divālā} again
\end{itemize}

\begin{tcolorbox}[breakable,colback=teal!4,colframe=teal!35!black,title={Before you move on}]
What does \hi{दिवाला} mean? (\textquotedblleft{}Bankruptcy\textquotedblright{}.) Say it once more.
\end{tcolorbox}
\section[\texorpdfstring{\hi{निधि} (nidhi)}{निधि (nidhi)}]{\hi{निधि} (nidhi) — a fund; a treasure}

\label{lesson:HI-C430-nidhi}



\begin{quote}
Before the new one: say the Hindi for to cash, then the Hindi for bankruptcy.
\end{quote}

\subsection*{You'll want to know: \hi{निधि}}
\textbf{\hi{निधि}} — \emph{nidhi} — \textquotedblleft{}a fund; a treasure\textquotedblright{}.

\begin{quote}
\hi{गाँव} \hi{के} \hi{लोगों} \hi{ने} \hi{स्कूल} \hi{के} \hi{लिए} \hi{एक} \hi{निधि} \hi{बनाई।}
\end{quote}

\begin{quote}
\emph{gā̃v ke logoṁ ne skūl ke liye ek nidhi banāī.}
\end{quote}

\begin{quote}
The villagers set up a fund for the school.
\end{quote}

Feminine.

\begin{cousinweb}
\textbf{\hi{निधि}} (\emph{nidhi}) is a tatsama, Sanskrit \emph{nidhi}, \textquotedblleft{}a place for storing; a store, a treasure\textquotedblright{} (Monier-Williams); banks and charities use it for a fund.
\end{cousinweb}

\begin{grammarlens}[title={What you've built}]
One more word for getting around.
\end{grammarlens}

\subsection*{Your turn}
\begin{itemize}
\raggedright
  \item \emph{Say it:} \emph{nidhi}
  \item \emph{Say it:} \emph{nidhi}, once more
  \item \emph{Say it:} \emph{gā̃v ke logoṁ ne skūl ke liye ek nidhi banāī.} — The villagers set up a fund for the school.
  \item \emph{Recall:} say the Hindi for to cash, then the Hindi for bankruptcy, then say \emph{nidhi} again
\end{itemize}

\begin{tcolorbox}[breakable,colback=teal!4,colframe=teal!35!black,title={Before you move on}]
What does \hi{निधि} mean? (\textquotedblleft{}A fund; a treasure\textquotedblright{}.) Say it once more.
\end{tcolorbox}
\section[\texorpdfstring{\hi{पूँजी} (pū̃jī)}{पूँजी (pū̃jī)}]{\hi{पूँजी} (pū̃jī) — capital (money); savings}

\label{lesson:HI-C430-puji}



\begin{quote}
Before the new one: say the Hindi for bankruptcy, then the Hindi for a fund.
\end{quote}

\subsection*{You'll want to know: \hi{पूँजी}}
\textbf{\hi{पूँजी}} — \emph{pū̃jī} — \textquotedblleft{}capital (money); savings\textquotedblright{}.

\begin{quote}
\hi{दुकान} \hi{खोलने} \hi{के} \hi{लिए} \hi{पूँजी} \hi{चाहिए।}
\end{quote}

\begin{quote}
\emph{dukān kholne ke liye pū̃jī cāhiye.}
\end{quote}

\begin{quote}
You need capital to open a shop.
\end{quote}

Feminine.

\begin{cousinweb}
\textbf{\hi{पूँजी}} (\emph{pū̃jī}) is a tadbhava of Sanskrit \emph{puñja}, \textquotedblleft{}a heap, a quantity\textquotedblright{} (CDIAL 8251), which lists Hindi \emph{pū̃jī}, \textquotedblleft{}heap, store\textquotedblright{}: the store of money a business is built on.
\end{cousinweb}

\begin{grammarlens}[title={What you've built}]
One more word for getting around.
\end{grammarlens}

\subsection*{Your turn}
\begin{itemize}
\raggedright
  \item \emph{Say it:} \emph{pū̃jī}
  \item \emph{Say it:} \emph{pū̃jī}, once more
  \item \emph{Say it:} \emph{dukān kholne ke liye pū̃jī cāhiye.} — You need capital to open a shop.
  \item \emph{Recall:} say the Hindi for bankruptcy, then the Hindi for a fund, then say \emph{pū̃jī} again
\end{itemize}

\begin{tcolorbox}[breakable,colback=teal!4,colframe=teal!35!black,title={Before you move on}]
What does \hi{पूँजी} mean? (\textquotedblleft{}Capital (money); savings\textquotedblright{}.) Say it once more.
\end{tcolorbox}
\section[\texorpdfstring{\hi{चुकना} (cuknā)}{चुकना (cuknā)}]{\hi{चुकना} (cuknā) — to be paid off; to run out}

\label{lesson:HI-C430-cuknaintr}



\begin{quote}
Before the new one: say the Hindi for a fund, then the Hindi for capital (money).
\end{quote}

\subsection*{You'll want to know: \hi{चुकना}}
\textbf{\hi{चुकना}} — \emph{cuknā} — \textquotedblleft{}to be paid off; to run out\textquotedblright{}.

\begin{quote}
\hi{हमारा} \hi{क़र्ज़} \hi{अब} \hi{चुक} \hi{गया} \hi{है।}
\end{quote}

\begin{quote}
\emph{hamārā qarz ab cuk gayā hai.}
\end{quote}

\begin{quote}
Our loan is paid off now.
\end{quote}

Past (it …): \hi{वह} \hi{क़र्ज़} \hi{चुक} \hi{गया}, \emph{vah qarz cuk gayā}.

\begin{cousinweb}
\textbf{\hi{चुकना}} (\emph{cuknā}) is the intransitive partner of \hi{चुकाना}, from \emph{*cukk}, \textquotedblleft{}fall short of, stop\textquotedblright{}, a form CDIAL reconstructs (CDIAL 4848), which lists Hindi \emph{cuknā}, \textquotedblleft{}to be finished\textquotedblright{}.
\end{cousinweb}

\begin{grammarlens}[title={What you've built}]
One more word for getting around.
\end{grammarlens}

\subsection*{Your turn}
\begin{itemize}
\raggedright
  \item \emph{Say it:} \emph{cuknā}
  \item \emph{Say it:} \emph{cuknā}, once more
  \item \emph{Say it:} \emph{hamārā qarz ab cuk gayā hai.} — Our loan is paid off now.
  \item \emph{Recall:} say the Hindi for a fund, then the Hindi for capital (money), then say \emph{cuknā} again
\end{itemize}

\begin{tcolorbox}[breakable,colback=teal!4,colframe=teal!35!black,title={Before you move on}]
What does \hi{चुकना} mean? (\textquotedblleft{}To be paid off; to run out\textquotedblright{}.) Say it once more.
\end{tcolorbox}
\section[\texorpdfstring{\hi{उगाहना} (ugāhnā)}{उगाहना (ugāhnā)}]{\hi{उगाहना} (ugāhnā) — to collect (dues); to raise (donations)}

\label{lesson:HI-C430-ugahna}



\begin{quote}
Before the new one: say the Hindi for capital (money), then the Hindi for to be paid off.
\end{quote}

\subsection*{You'll want to know: \hi{उगाहना}}
\textbf{\hi{उगाहना}} — \emph{ugāhnā} — \textquotedblleft{}to collect (dues); to raise (donations)\textquotedblright{}.

\begin{quote}
\hi{हमने} \hi{मेले} \hi{के} \hi{लिए} \hi{मोहल्ले} \hi{से} \hi{पैसे} \hi{उगाहे।}
\end{quote}

\begin{quote}
\emph{hamne mele ke liye mohalle se paise ugāhe.}
\end{quote}

\begin{quote}
We collected money from the neighbourhood for the fair.
\end{quote}

Past, with \hi{ने} (he or she …): \hi{उसने} \hi{किराया} \hi{उगाहा}, \emph{usne kirāyā ugāhā}.

\begin{cousinweb}
\textbf{\hi{उगाहना}} (\emph{ugāhnā}) is a tadbhava of Sanskrit \emph{udgrāhayati}, \textquotedblleft{}causes to take up, causes to pay\textquotedblright{} (CDIAL 1967), which lists Hindi \emph{ugāhnā}, \textquotedblleft{}to accumulate\textquotedblright{}.
\end{cousinweb}

\begin{grammarlens}[title={What you've built}]
One more word for getting around.
\end{grammarlens}

\subsection*{Your turn}
\begin{itemize}
\raggedright
  \item \emph{Say it:} \emph{ugāhnā}
  \item \emph{Say it:} \emph{ugāhnā}, once more
  \item \emph{Say it:} \emph{hamne mele ke liye mohalle se paise ugāhe.} — We collected money from the neighbourhood for the fair.
  \item \emph{Recall:} say the Hindi for capital (money), then the Hindi for to be paid off, then say \emph{ugāhnā} again
\end{itemize}

\begin{tcolorbox}[breakable,colback=teal!4,colframe=teal!35!black,title={Before you move on}]
What does \hi{उगाहना} mean? (\textquotedblleft{}To collect (dues); to raise (donations)\textquotedblright{}.) Say it once more.
\end{tcolorbox}
